% !TEX TS-program = pdflatex
% English Grammar Comprehensive Guide - Mosioc
% Comprehensive LaTeX reference book aimed at IELTS & TOEFL learners

\documentclass[12pt,oneside]{book}

% Packages for layout and typography
\usepackage[utf8]{inputenc}
\usepackage[T1]{fontenc}
\usepackage{lmodern}
\usepackage{microtype}
\usepackage{geometry}
\geometry{a4paper, left=30mm, right=25mm, top=30mm, bottom=30mm}
\usepackage{setspace}
\onehalfspacing
\usepackage{titlesec}
\usepackage{xcolor}
\usepackage{tcolorbox}
\tcbuselibrary{skins, breakable}
\usepackage{hyperref}
\hypersetup{colorlinks=true, linkcolor=blue, urlcolor=blue}
\usepackage{enumitem}
\usepackage{amsmath,amssymb}
\usepackage{booktabs}
\usepackage{longtable}
\usepackage{multicol}
\usepackage{caption}
\usepackage{graphicx}
\usepackage{framed}
\usepackage{fancyhdr}
\usepackage{listings}
\usepackage{imakeidx}
\makeindex

% Styling for chapter and section headings
\titleformat{\chapter}[display]
  {\normalfont\huge\bfseries}{\chaptertitlename\ \thechapter}{20pt}{\Huge}
\titleformat{\section}{\large\bfseries}{\thesection}{1em}{}
\titleformat{\subsection}{\normalsize\bfseries}{\thesubsection}{1em}{}

% Boxes for notes, tips, examples
\newtcolorbox{examptip}{colback=blue!5!white,colframe=blue!50!black,breakable}
\newtcolorbox{examplebox}{colback=green!5!white,colframe=green!50!black,breakable}
\newtcolorbox{commonerr}{colback=red!5!white,colframe=red!60!black,breakable}
\newtcolorbox{definitionbox}{colback=gray!5!white,colframe=black!30!white,breakable}

% Commands
\newcommand{\ex}[1]{\begin{examplebox}\textbf{Example:} #1\end{examplebox}}
\newcommand{\tip}[1]{\begin{examptip}\textbf{Exam Tip:} #1\end{examptip}}
\newcommand{\cerror}[1]{\begin{commonerr}\textbf{Common Error:} #1\end{commonerr}}
\newcommand{\defbox}[1]{\begin{definitionbox}\textbf{Definition:} #1\end{definitionbox}}

% Document metadata
\title{English Grammar: A Comprehensive Reference for IELTS \& TOEFL}
\author{Mosioc}
\date{\today}

\begin{document}

% Title page
\begin{titlepage}
  \centering
  \vspace*{3cm}
  {\Huge\bfseries English Grammar}\vspace{0.5cm} \\
  {\Large\bfseries A Comprehensive Reference for IELTS \& TOEFL}\\[1.5cm]
  {\Large Mosioc}\vfill
  {\large Compiled and Expanded — \today}\\[1cm]
  \vfill
  \normalsize This textbook is designed as a practical and in-depth grammar reference for English learners preparing for academic exams (IELTS \& TOEFL). It focuses on both explanation and application: rules, examples, exam tips, and practice.
\end{titlepage}

\frontmatter
\tableofcontents
\mainmatter

% Preface
\chapter*{Preface}
This book is written for intermediate to advanced English learners who aim to achieve high scores in standardized academic tests such as \textbf{IELTS} and \textbf{TOEFL}, and for learners who want a comprehensive, reliable grammar reference. Each chapter includes clear explanations, graded examples, exam tips, and practice exercises.

\tip{Use this book actively: read a section, study the examples, then do the exercises. When preparing for IELTS/TOEFL, focus on accuracy first, range second.}

\chapter{Foundations of English Grammar}
\section{What is Grammar?}
Grammar is the system of rules that governs how words are combined to form sentences. Grammar controls morphology (word forms), syntax (sentence structure), and usage (pragmatic choices).

\section{Sentence Elements}
\defbox{A sentence typically contains a \textbf{subject} and a \textbf{predicate}. The predicate contains the verb and provides information about the subject. Objects, complements, and modifiers add details.}

\subsection{Subject, Predicate, Object, Complement}
\begin{itemize}
  \item \textbf{Subject:} the doer or topic of the sentence.\
  \item \textbf{Predicate:} includes the verb and expresses something about the subject.\
  \item \textbf{Direct object:} receives the action of a transitive verb.\
  \item \textbf{Indirect object:} recipient of the direct object.\
  \item \textbf{Complement:} completes the meaning of a subject or object (e.g. subject complement after linking verbs).\
\end{itemize}

\ex{\textit{The committee} (subject) \textit{approved} (verb) \textit{the proposal} (direct object).}

\section{Word Order in English}
English is primarily SVO (Subject-Verb-Object). Word order is crucial for meaning. In questions and some other constructions, inversion occurs.

\chapter{Parts of Speech}
\section{Nouns}
\subsection{Types of Nouns}
Common, proper, countable, uncountable, collective, abstract.

\subsection{Countable vs. Uncountable}
\defbox{Countable nouns have plural forms and can be used with numbers (one book, two books). Uncountable nouns cannot be counted directly (water, information) and take singular verbs.}

\ex{\textit{Many books} vs. \textit{much information}.}
\cerror{Using 'many' with uncountable nouns: *many informations* is incorrect.}

\section{Pronouns}
\subsection{Personal, Possessive, Reflexive}
Rules for pronoun agreement and clarity of reference.

\ex{\textit{She gave herself a break.}} 
\cerror{Ambiguous pronouns: avoid unclear antecedents.}

\section{Adjectives and Adverbs}
\subsection{Adjective order}
Typical order: Determiner — Opinion — Size — Age — Shape — Color — Origin — Material — Purpose — Noun.\
\ex{\textit{a lovely small old round wooden jewelry box}}

\subsection{Adverbs: Position and Types}
Adverbs of frequency, manner, place, time; typical positions (before main verb for frequency, after main verb for manner, etc.).

\chapter{Sentence Structure and Clauses}
\section{Simple, Compound, Complex}
\defbox{\textbf{Simple sentence:} one independent clause. \\ \textbf{Compound sentence:} two or more independent clauses joined by coordinating conjunctions (and, but, or, so, yet). \\ \textbf{Complex sentence:} one independent clause + one or more dependent (subordinate) clauses.}

\ex{Simple: \textit{She reads.} Compound: \textit{She reads, and he writes.} Complex: \textit{She reads because she enjoys learning.}}

\section{Relative Clauses}
Defining vs non-defining clauses; punctuation rules, choice of relative pronoun (who, which, that, whom, where, when).

\ex{The book \textit{that} I borrowed was helpful. (defining) \\ My sister, \textit{who} lives in London, is a teacher. (non-defining)}

\section{Noun Clauses (Nominal Clauses)}
Noun clauses function as subjects, objects, or complements. Common introducers: that, whether, if, wh-words.

\ex{\textit{What she said} surprised everyone.}

\section{Sentence Fragments and Run-ons}
How to identify and correct fragments and comma splices (run-on sentences). Use semicolons, coordinating conjunctions, or restructure clauses.

\chapter{Tenses and Aspect}
\section{Understanding Tense vs Aspect}
Tense locates an action in time (past, present, future). Aspect shows internal temporal structure (simple, continuous/progressive, perfect, perfect continuous).

\section{Present Tenses}
\subsection{Present Simple}
Use for habitual actions, general truths, scheduled events.
\ex{She usually walks to work.}
\subsection{Present Continuous}
Use for actions happening now or around now, and future arrangements.
\ex{I am studying for the exam.}
\subsection{Present Perfect}
Used for actions with relevance to the present: experience, recent events, actions that started in the past and continue to present.
\ex{She has lived here since 2010.}
\subsection{Present Perfect Continuous}
Focuses on duration of an action up to the present.
\ex{They have been studying for three hours.}

\section{Past Tenses}
\subsection{Past Simple}
Completed actions at a definite time in the past.
\ex{He graduated last year.}
\subsection{Past Continuous}
Background actions in the past; parallel actions.
\ex{While I was cooking, she called.}
\subsection{Past Perfect}
An action completed before another past action.
\ex{She had left before I arrived.}
\subsection{Past Perfect Continuous}
Duration before a certain past point.
\ex{They had been working for two hours when the power went out.}

\section{Future Forms}
Will vs going to vs present continuous for future. Future perfect and future continuous explained.

\section{Tense Harmony and Sequence of Tenses}
How tense backshifts in reported speech and conditionals. Importance of consistency for coherence in essays and spoken answers.

\tip{In IELTS Writing Task 1 (Academic), prefer simple and perfect tenses to describe trends. Avoid unnecessary tense shifts.}

\chapter{Modal Verbs and Auxiliary Structures}
\section{Modal Verbs: basic uses}
Can/could, may/might, must, should, ought to, shall/will, would.

\subsection{Degrees of certainty and politeness}
Can vs could, may vs might, must (deduction) vs have to (obligation).

\ex{He must be at home. (logical conclusion) \\ You must finish this by tomorrow. (obligation)}

\section{Semi-modals and Periphrastic forms}
Have to, need to, be able to, used to, be supposed to.

\chapter{Voice: Active and Passive}
\section{Forming the Passive}
All tenses can have passive equivalents. Use auxiliary \textit{be} + past participle.

\ex{Active: The committee approved the policy. \\ Passive: The policy was approved by the committee.}

\section{When to use Passive}
To focus on action rather than agent, in academic writing, or when agent is unknown or irrelevant.

\tip{In IELTS Academic Writing Task 1, passive voice can often make descriptions more objective.}

\chapter{Conditionals and Hypothetical Structures}
\section{Zero, First, Second, Third Conditionals}
\begin{itemize}
  \item Zero: general truths (If + present, present). \
  \item First: possible future (If + present, will + base). \
  \item Second: unreal present/future (If + past simple, would + base). \
  \item Third: unreal past (If + past perfect, would have + past participle). \
\end{itemize}

\ex{If you heat water, it boils. \\ If it rains, I will stay home. \\ If I won the lottery, I would travel. \\ If I had known, I would have told you.}

\section{Mixed Conditionals and Wishes}
Mixed conditionals combine past and present hypothetical consequences. Use \textit{wish} to express regret or desire with past/present/future reference.

\ex{If I had studied harder (past), I would be at university now (present result). \\ I wish I had gone to the conference.}

\chapter{Clauses and Advanced Structures}
\section{Relative Clauses: reduction and complexity}
Reduced relative clauses, omission of relative pronouns when possible, and restrictive vs nonrestrictive distinctions.

\section{Inversion for Emphasis}
Inversion in negatives and for adverbial expressions for stylistic emphasis.

\ex{Rarely have I seen such dedication.}

\section{Cleft and Pseudo-cleft Sentences}
It was/what-clefts to emphasize a sentence element. Useful to vary sentence openings in essays.

\ex{It was \textbf{John} who solved the problem. \\ What I need is more time.}

\section{Nominalization}
Turning verbs and adjectives into nouns to create formal, concise academic style (e.g. \textit{analyze} -> \textit{analysis}).

\tip{Nominalization increases formality and density—use it in academic writing but avoid overuse.}

\chapter{Reported Speech}
\section{Backshifting and Reporting Verbs}
Rules for shifting tenses, pronouns, and time expressions when converting direct speech into reported speech.

\ex{Direct: "I am leaving" -> Reported: He said (that) he was leaving.}

\chapter{Subjunctive and Formal Structures}
\section{Present Subjunctive}
Use of subjunctive in formal expressions, demands, suggestions (e.g. "It is essential that he be present").

\section{Formal registers and fixed expressions}
Use of constructions like \textit{lest}, \textit{had I known}, and formal inversion.

\chapter{Punctuation, Style and Cohesion}
\section{Articles and Determiners}
Usage of a/an, the, zero article; determiners with countable and uncountable nouns.

\ex{I bought a car. The car is blue.}

\section{Commas, Semicolons, Colons}
Rules to avoid comma splices and to use semicolons and colons effectively.

\section{Linking words and discourse markers}
Phrases for addition, contrast, cause/effect, exemplification, sequencing. Use them to improve cohesion and coherence in essays.

\begin{multicols}{2}
\textbf{Addition:} moreover, furthermore, in addition, besides.\\
\textbf{Contrast:} however, on the other hand, whereas, although.\\
\textbf{Cause/Effect:} therefore, consequently, as a result.\\
\textbf{Example:} for example, for instance, such as.\\
\textbf{Sequencing:} firstly, secondly, finally, subsequently.
\end{multicols}

\tip{Avoid overusing the same linking word. Aim for a varied set of connectors and natural flow.}

\chapter{Common Errors and How to Avoid Them}
\section{Subject-Verb Agreement}
Agreement with singular/plural subjects, collective nouns, and tricky constructions (e.g., "one of the students is/are").

\section{Article Errors}
When to use definite/indefinite/zero articles; common mistakes by L1 speakers.

\section{Preposition Problems}
Common preposition collocations and fixed expressions. Practice lists.

\cerror{Incorrect preposition: *married with* -> correct: *married to*.}

\chapter{Grammar for IELTS and TOEFL}
\section{What examiners look for}
IELTS: Grammatical Range and Accuracy. TOEFL: grammatical control in writing and speaking responses.

\section{Writing Task 1 (Academic) and Writing Task 2}
Task-specific grammar guidance: report trends, use passive where appropriate, vary sentence types, avoid contractions.

\section{Speaking: Fluency and Accuracy}
Balance fluency with grammatical correctness. Use a range of structures but avoid unnatural complexity.

\tip{When practicing speaking, record yourself and transcribe the response to analyze grammar and sentence variety.}

\chapter{Exercises}
Each section below contains short practice activities.

\section{Error Correction}
Correct the errors in the following sentences:
\begin{enumerate}
  \item She don't like coffee.
  \item The informations were useful.
  \item If I will see him, I will tell him.
  \item He suggested me to go there.
  \item This is the man which helped me.
\end{enumerate}

\section{Transformations}
Rewrite the sentences as requested:
\begin{enumerate}
  \item Change to passive: "They will announce the results tomorrow." 
  \item Reported speech: "She said, 'I have finished the report.'"
  \item Conditional (second): "If I (win) the prize, I (buy) a car." 
  \item Use a cleft to emphasize "John fixed the car." 
  \item Reduce the relative clause: "The documents which were signed yesterday are ready." 
\end{enumerate}

\section{Tense Practice}
Fill in the correct tense:
\begin{enumerate}
  \item By next year, I (finish) my degree.
  \item She (study) French for three years before she moved to Paris.
  \item Right now, they (work) on the new project.
\end{enumerate}

\chapter*{Answer Key}
\addcontentsline{toc}{chapter}{Answer Key}
\section*{Error Correction Answers}
\begin{enumerate}
  \item She doesn't like coffee.
  \item The information was useful. (information is uncountable)
  \item If I see him, I will tell him. (first conditional) or If I saw him, I would tell him. (second conditional)
  \item He suggested that I go there. / He suggested going there.
  \item This is the man who helped me. (use "who")
\end{enumerate}

\section*{Transformations Answers}
\begin{enumerate}
  \item The results will be announced tomorrow.
  \item She said (that) she had finished the report.
  \item If I won the prize, I would buy a car.
  \item It was John who fixed the car.
  \item The documents signed yesterday are ready.
\end{enumerate}

\section*{Tense Practice Answers}
\begin{enumerate}
  \item will have finished
  \item had been studying
  \item are working / are working on
\end{enumerate}

\chapter*{Appendix: Useful Tables}
\addcontentsline{toc}{chapter}{Appendix: Useful Tables}
\section*{Timeline of Tenses}
\begin{longtable}{p{4cm}p{10cm}}
\toprule
Tense & Use / Example \\
\midrule
Present Simple & habits, facts: She works every day. \\
Present Continuous & actions now, arrangements: He is leaving tomorrow. \\
Present Perfect & experience/continuing: She has lived here since 2010. \\
Past Simple & finished past: They visited last summer. \\
Past Perfect & past before past: She had left when I arrived. \\
Future (will) & predictions/decisions: I will help you. \\
\bottomrule
\end{longtable}

\printindex

\end{document}
